% Source: https://artofproblemsolving.com/community/c3h2244020_lcm_fractions
% Stable identity: aops-topic-2244020-q01
Write each number in lowest terms $\dfrac {a}{b}$, and a candidate as $\dfrac {p}{q}$ in lowest terms.

\textbf{(a)} $\dfrac {p}{q}\div\dfrac {a}{b}=\dfrac{pb}{qa}$ is an integer exactly when $a\mid p$ and $q\mid b$: if $qa\mid pb$, then $q\mid pb$ with $\gcd(q,p)=1$ gives $q\mid b$, and $a\mid pb$ with $\gcd(a,b)=1$ gives $a\mid p$; the converse is clear. So $\dfrac {p}{q}$ is a common multiple exactly when $p$ is a multiple of every numerator and $q$ divides every denominator. The smallest choice is
\[\dfrac{\operatorname{lcm}(2,9,6,3)}{\gcd(5,4,7,8)}=\dfrac{18}{1}=\boxed{18}.\]
Check: $18$ divided by the four numbers gives $45$, $8$, $21$, $48$.

\textbf{(b)} In the same way $\dfrac {a}{b}\div\dfrac {p}{q}=\dfrac{aq}{bp}$ is an integer exactly when $p\mid a$ and $b\mid q$. The largest common divisor is
\[\dfrac{\gcd(2,9,6,3)}{\operatorname{lcm}(5,4,7,8)}=\boxed{\dfrac{1}{280}}.\]
Check: the four numbers divided by $\dfrac{1}{280}$ give $112$, $630$, $240$, $105$.
